\section{Compartición de conocimiento e iteraciones futuras}

\subsection{Flujo de conocimiento y reproducción}
El equipo de AlphaProof registra el experiment log, el failed proof trace, los human comment y los experimentos del newsletter en un solo repo; cada release viene acompañado de un `diff` que se usa para reentrenar la policy. Hubert señala: «Necesitamos compartir con los nuevos integrantes los aspectos destacados de la QA, las governance decision y los proof de ejemplos positivos y negativos, para que el training loop pueda iterar más rápido».
\begin{importantbox}{Mecanismo de compartición de conocimiento}
1) Experiment repo, que incluye el chunk gating journal; 2) Newsletter, que registra cada semana los aspectos destacados de QA + governance; 3) Shared slides + transcript, que también son el material para el ramp-up de los nuevos integrantes.
\end{importantbox}

\subsection{Aspectos destacados de la QA}
La sesión de QA se usó para explicar si AlphaProof puede extenderse a la non-Euclidean geometry y al ~hypergraph theorem. Al responder, Hubert subrayó que la RL policy debe diseñarse de forma modular, ajustando cada vez solo el search depth o el tactic mix, para poder extenderse gradualmente a demostraciones más abstractas.
\begin{warningbox}{Puntos de precaución en la QA}
No se debe usar AlphaProof de manera indiscriminada para resolver todos los problemas; en los idiomas low-resource o en las novel math branch, la policy necesita un nuevo warm start, junto con el chunk gating y el fairness board.
\end{warningbox}

\subsection{Resumen de la sección}
El apartado sobre compartición de conocimiento mostró cómo el log, el newsletter y el transcript forman una red de reproducción, y cómo los retos de extensión expuestos mediante la sesión de QA explican mejor los límites de la iteration y las necesidades de gobernanza.
\newpage

